\section{Optical moments below the homogeneous endpoint}
\label{sec:v88-moments}
The unconditioned incidence estimate permits a larger moment range than
its joint endpoint-localized counterpart. We keep these two estimates
separate. In particular, no improvement of the mixed inverse-gradient
moment is inferred from the theorem below.

Fix one admissible finite roof-frequency band $B_\dagger$. All constants
in this section may depend on this band and on the displayed exponents.
The band never depends on the collision count, an incidence threshold,
or a physical roof width.

\subsection{The retained spectral assumptions}
Write $Q=\mathscr T_{R,\xi}$, where $\xi$ is the real four-record
frequency, $M_s=M_{\one_{\{c<s\}}}$, and $\ell$ is mass evaluation.
We use the norm, multiplier, and physical-realization statements of
Supplement module 183, together with its fixed-band spectral input.
Specifically, the initial vector is $\one$, its physical iterates satisfy
$|Q^j\one|=1$, and
\begin{align}
 \|M_s f\|&\le C\|f\|,& |M_s f|_w&\le Cs^{1/8}\|f\|,
                                                    \label{eq:v88-multiplier}\\
 |M_sQ^j\one|_w&\le Cs^{3/2},&\|Q^j\|&\le C.
                                                    \label{eq:v88-real-weak}
\end{align}
The second estimate in \eqref{eq:v88-real-weak} is used on peak
charts. The first follows from homogeneous stable-curve length, in the
original invariant-density convention: incidence is $O(k^{-2})$ and
a homogeneous stable curve has length $O(k^{-3})$. No additional cosine
is inserted in the weak norm. Physical strong representatives and their
weak limits are the ones specified in the retained multiplier theorem.

On each peak chart the retained decomposition and inequalities are
\begin{align}
 Q^j&=\lambda^j\Pi+N^j,&\|N^j\|&\le C\rho^j,\notag\\
 \|Q^j f\|&\le C\sigma^j\|f\|+C(1+j)|f|_w,&
 |Q^j f|_w&\le C(1+j)|f|_w,                   \label{eq:v88-spectral}
\end{align}
with $\rho,\sigma<1$. The fixed-band complementary region has a
uniform spectral gap. The real peak powers have the four-dimensional
Gaussian integral bound $O(m^{-2})$. This paragraph states inherited
four-record inputs; it does not attribute them to a theorem for a
smaller record.

\begin{lemma}[Adjustable weak interpolation]\label{lem:v88-interpolation}
Let $0<\theta<1$ and $0<\epsilon_0<3/2$. A fixed refinement of the
peak charts gives
\begin{align}
 \|\Pi f\|+\sup_{a\ge0}|\ell(Q^a f)|
       &\le C\|f\|^{1-\theta}|f|_w^\theta,
                                        \label{eq:v88-interpolation}\\
 |M_s\Pi\one|_w&\le Cs^{3/2-\epsilon_0}.
                                        \label{eq:v88-projected}
\end{align}
All unit-modulus loci are retained. No normalization by $\Pi\one$ is
used, and the lemma remains valid when that vector vanishes.
\end{lemma}
\begin{proof}
Choose $\gamma>0$ with $1-\gamma>\theta$ and shrink the charts so that
$|\lambda|\ge r_*$ and
\[
 \frac{\log(1/r_*)}{\log(1/\sigma)}<\gamma,\qquad
 \frac{3\log(1/r_*)}{2\log(1/\rho)}<\epsilon_0.
\]
Uniform constants may first be chosen on larger fixed neighborhoods.
The remaining compact portions have a spectral gap and are still
estimated on the complement; they are not removed from the record or
from the final arithmetic reference.

Normalize $\|f\|=1$ and let $t=|f|_w<1/2$. Choose $L$ with
$\sigma^L\le t<\sigma^{L-1}$. Commutation and
\eqref{eq:v88-spectral} give
$\|\Pi f\|\le C(1+|\log t|)t^{1-\gamma}\le Ct^\theta$.
For $a\le L$ use the weak power inequality and continuity of mass;
for $a>L$ use the strong inequality at $L$ and bounded powers
thereafter. The cases $t=0$ and $t$ bounded below follow by passage to the limit
and boundedness, respectively. Scaling proves the first assertion.

For the second choose
$L=\lceil(3/2)\log(1/s)/\log(1/\rho)\rceil$ in the identity
\[
 M_s\Pi\one=\lambda^{-L}M_sQ^L\one-
                         \lambda^{-L}M_sN^L\one.
\]
Its weak norm is at most
$Cr_*^{-L}(s^{3/2}+s^{1/8}\rho^L)\le Cs^{3/2-\epsilon_0}$.
Its strong norm is bounded by \eqref{eq:v88-multiplier}. This uses
physical iterates, not an asserted modulus-one representative of a
continued eigenvector. The auxiliary integer $L$ is an interpolation
length, not a Fourier-band choice.
\end{proof}

\begin{theorem}[A capped marked estimate below $3/2$]
\label{thm:v88-incidence}
For every $0<\kappa<3/2$ and the fixed band $B_\dagger$, there
are $C_{\kappa,B_\dagger}<\infty$ and $\rho_\dagger<1$ such that
\begin{align}
 &m^2\nu\{K_m^{\rm c}=k,A_m=n,F\in J,c(T_R^j x)<s\}\notag\\
 &\quad\le C_{\kappa,B_\dagger}s^\kappa(|J|+B_\dagger^{-1})
  +C_{\kappa,B_\dagger}(1+B_\dagger|J|)m^2
                      \min\{\rho_\dagger^{m/2},s^2\}.
                                             \label{eq:v88-incidence}
\end{align}
The estimate holds for every finite $m\ge2$, every mark $0\le j\le m$,
interval $J$, exact coefficient and sufficiently small $s>0$.
No assertion is made at $\kappa=3/2$.
\end{theorem}
\begin{proof}
Choose $\theta,\epsilon_0$ so that
$(3/2-\epsilon_0)\theta>\kappa$. In the exact pairing
$\ell(Q^{m-j}M_sQ^j\one)$, at least one chronological block is
of length $m/2$. If the left block is long, expand it and apply
\eqref{eq:v88-interpolation} to $M_sQ^j\one$. The amplitude is
$O(s^{3\theta/2})$. If the right block is long, expand it and
apply \eqref{eq:v88-interpolation} to $M_s\Pi\one$, using
\eqref{eq:v88-projected}. The amplitude is
$O(s^{(3/2-\epsilon_0)\theta})$. Both are $O(s^\kappa)$.
The long-block remainders and the complementary frequency region are
$O(\rho_\dagger^{m/2})$, after changing fixed constants.

At every real frequency the same physical pairing also has absolute
value at most $\nu\{c<s\}=1-\sqrt{1-s^2}\le s^2$.
For nonnegative $A,D,V$, the two bounds $|z|\le A+D$ and
$|z|\le V$ imply $|z|\le A+\min(D,V)$. Thus the entire remainder
retains this physical cap. Apply the positive interval majorant with
roof-frequency support in the fixed band. Its Fourier supremum is
$O(|J|+B_\dagger^{-1})$. Gaussian integration of each long-block
peak term gives $m^{-2}$, whereas integration of the capped remainder
gives the second term of \eqref{eq:v88-incidence}. Exact torus
coefficient extraction retains $k,n$ separately. The finite refinement
has only changed the constants at the fixed band.
\end{proof}

\begin{theorem}[Actual complete optical moments below $3/2$]
\label{thm:v88-moment}
For every fixed $1<p<3/2$, the actual density
$u_p(t)\,dt=F_\#(I_m^p\one_\Pi\nu_R^*)$ is finite at every
collision count. On the one fixed band, for all intervals $J$,
\begin{align}
 m^2\int_Ju_p(t)\,dt\le C_{p,B_\dagger}(|J|+B_\dagger^{-1})
 +C_{p,B_\dagger}(1+B_\dagger|J|)m^2
                    \rho_\dagger^{m(2-p)/4}. \label{eq:v88-moment}
\end{align}
In particular, for $|J|\le1$,
\begin{equation}\label{eq:v88-uniform-energy}
       \sup_{m,R,n,k}m^2\int_Ju_p\le K_{p,B_\dagger}<\infty.
\end{equation}
No singular multiplier on the strong space is asserted. These are
unconditioned-gradient moments, not mixed $I_m^pz^{-\eta}$ bounds.
\end{theorem}
\begin{proof}
Choose $p<\kappa<3/2$. The positive incidence layer cake applied to
\eqref{eq:v88-incidence} has leading integral
$\int_0^{s_0}s^{\kappa-p-1}\,ds<\infty$. For the full remainder,
\[
 s_0^{-p}\min(D,s_0^2)+p\int_0^{s_0}s^{-p-1}\min(D,s^2)\,ds
      =\frac2{2-p}\min\{s_0^{2-p},D^{1-p/2}\}.
\]
This follows by splitting at $\sqrt D$ and holds for $p<2$.
Use $D=\rho_\dagger^{m/2}$ before removing any threshold. The
bounded-incidence part uses the retained unweighted fixed-band
concentration estimate.

Let $S=1+2q_0/(1-q_0)$. Weighted Jensen gives
$I_m^p\le S^{p-1}(1+\sum_iq_0^{d_i}c_i^{-p})$.
Sum the marked estimates with these geometric weights. Dropping the
endpoint tests on the positive upper comparison costs only $c_*^{-1}$
and leaves the exact coefficient unchanged; it does not redefine
$u_p$. This proves \eqref{eq:v88-moment}. The exponential remainder
has bounded $m^2$ supremum, proving \eqref{eq:v88-uniform-energy}.
Taking a containing roof interval proves finite mass at fixed count;
coarea on compact regular charts and positive exhaustion prove absolute
continuity in the common fiber version. Neither limit $p\uparrow3/2$
nor $B_\dagger\to\infty$ is used in the microscopic arguments below.
\end{proof}
